\documentclass[a4paper,12pt,oneside]{report}

% ==========================================
% KHAI BÁO CÁC GÓI HỖ TRỢ (PACKAGES)
% ==========================================
\usepackage[utf8]{inputenc}
\usepackage[T5]{fontenc}
\usepackage[vietnamese]{babel}
\usepackage{amsmath, amssymb} % Hỗ trợ công thức toán học
\usepackage{graphicx}         % Hỗ trợ chèn ảnh
\usepackage{hyperref}         % Hỗ trợ tạo link mục lục
\usepackage{geometry}         % Căn lề chuẩn luận văn
\geometry{
	a4paper,
	left=30mm,
	right=20mm,
	top=20mm,
	bottom=20mm,
}
\usepackage{titlesec}
\titleformat{\chapter}[display]{\normalfont\huge\bfseries\centering}{\chaptertitlename\ \thechapter}{20pt}{\Huge}

% ==========================================
% THÔNG TIN TRANG BÌA
% ==========================================
\title{
    \vspace{-2cm}
    \large ĐẠI HỌC QUỐC GIA TP.HCM \\
    TRƯỜNG ĐẠI HỌC KHOA HỌC TỰ NHIÊN \\
    KHOA CÔNG NGHỆ THÔNG TIN \\
    \vspace{3cm}
    \textbf{\huge KHÓA LUẬN TỐT NGHIỆP} \\
    \vspace{1.5cm}
    \textbf{\Large PHÁT TRIỂN HỆ THỐNG PHÂN ĐOẠN ẢNH X-QUANG VỀ XƯƠNG DỰA VÀO CẦU NHẮC TRỰC QUAN} \\
    \vspace{2cm}
}
\author{\textbf{Sinh viên thực hiện:} Thông Lúc}
\date{\vspace{4cm} TP. Hồ Chí Minh, Năm 2026}

% ==========================================
% BẮT ĐẦU TÀI LIỆU
% ==========================================
\begin{document}

% ==========================================
% CHƯƠNG 1: GIỚI THIỆU
% ==========================================
\chapter{GIỚI THIỆU}
\section{Bối cảnh chung}
\section{Động lực nghiên cứu}
\subsection{Động lực khoa học}
\subsection{Động lực thực tiễn}
\section{Phát biểu bài toán}
\subsection{Giai đoạn huấn luyện (Offline)}
\subsection{Giai đoạn suy luận (Online)}
\subsection{Cấu trúc ngữ nghĩa của câu nhắc}
\section{Thách thức của bài toán}
\section{Đóng góp của khóa luận}
\section{Bố cục của khóa luận}

% ==========================================
% CHƯƠNG 2: CÁC CÔNG TRÌNH LIÊN QUAN VÀ CƠ SỞ LÝ THUYẾT
% ==========================================
\chapter{CÁC CÔNG TRÌNH LIÊN QUAN VÀ CƠ SỞ LÝ THUYẾT}
\section{Các công trình liên quan}
\subsection{Các phương pháp phân đoạn ảnh y khoa tự động}
% Review về U-Net truyền thống và các biến thể phân đoạn xương mù (Blind Segmentation).
\subsection{Phương pháp học máy tương tác trong y khoa}
% Review về xu hướng Prompt-based learning (như Segment Anything Model - SAM) và sự cần thiết của Human-in-the-loop.

\section{Cơ sở lý thuyết nền tảng}
\subsection{Mạng Nơ-ron Tích chập (Convolutional Neural Networks)}
\subsection{Kiến trúc mạng U-Net và Cơ chế chú ý (Attention Mechanism)}

\section{Học máy dựa trên câu nhắc (Prompt-based Learning)}
\subsection{Biểu diễn không gian của câu nhắc trực quan}
% Cơ sở toán học để tạo Heatmap từ Bounding Box.
\subsection{Cơ chế chú ý có điều kiện (Conditioned Attention)}

\section{Phương pháp luận giải thích mô hình (Explainable AI - XAI)}
\subsection{Bản đồ nhiệt GradCAM}

\section{Các độ đo đánh giá hiệu năng (Evaluation Metrics)}
\subsection{Hệ số Dice và Intersection over Union (IoU)}
\subsection{Sai số đường biên HD95 (95\% Hausdorff Distance)}
\subsection{Độ chuẩn xác định vị tâm (Center-Based Localization - CBL)}

% ==========================================
% CHƯƠNG 3: MÔ HÌNH VÀ HỆ THỐNG ĐỀ XUẤT
% ==========================================
\chapter{MÔ HÌNH VÀ HỆ THỐNG ĐỀ XUẤT}
\section{Kiến trúc tổng thể của hệ thống}
\section{Tiền xử lý dữ liệu và loại bỏ nhiễu ảnh X-quang}
\section{Mô hình phân đoạn PGA-UNet (Prompt-Guided Attention U-Net)}
\subsection{Biểu diễn câu nhắc trực quan (Prompt Heatmap Representation)}
\subsection{Cổng không gian tại bộ mã hóa (Prompt Spatial Gate - Encoder)}
\subsection{Cơ chế chú ý có điều kiện tại bộ giải mã (Conditioned Attention Decoder)}
\section{Hệ thống phòng vệ và tinh chỉnh câu nhắc tự động (IPR Pipeline)}
\subsection{Cơ chế cảnh báo lỗi từ người dùng (Confidence \& Center Distance Check)}
\subsection{Giải pháp cứu hộ dự phòng bằng GradCAM}
\subsection{Thuật toán tinh chỉnh câu nhắc tuần tự (Iterative Prompt Refinement - IPR dựa trên Centroid)}
\section{Mô hình phân lớp bệnh lý bất thường}

% ==========================================
% CHƯƠNG 4: THỰC NGHIỆM VÀ ĐÁNH GIÁ KẾT QUẢ
% ==========================================
\chapter{THỰC NGHIỆM VÀ ĐÁNH GIÁ KẾT QUẢ}
\section{Thiết lập môi trường thực nghiệm}
\subsection{Tập dữ liệu ảnh X-quang (Dataset)}
\subsection{Cấu hình huấn luyện (Hyperparameters, Loss functions)}
\section{Đánh giá hiệu năng hệ thống phân đoạn (Phần lõi)}
\subsection{Đánh giá cơ sở (Baseline Comparison): U-Net vs. Attention U-Net vs. PGA-UNet}
\subsection{Đánh giá tính bền bỉ (Robustness): Kịch bản Zoom-out, Shift và Mixed}
\section{Đánh giá cơ chế phòng vệ và Tự tinh chỉnh (Ablation Study)}
\subsection{Hiệu năng của hệ thống gợi ý cứu hộ GradCAM}
\subsection{Sự hội tụ của IPR: Phân tích đánh đổi giữa số vòng lặp và độ chính xác}
\section{Trực quan hóa kết quả (Qualitative Results)}
\subsection{Trực quan hóa các kịch bản phân đoạn tiêu biểu}
\subsection{Trực quan hóa quá trình nắn nét tự động của IPR}
\section{Đánh giá kết quả mô hình phân lớp}
\section{Tích hợp và đánh giá ứng dụng thực tiễn (app.py)}

% ==========================================
% CHƯƠNG 5: KẾT LUẬN VÀ HƯỚNG PHÁT TRIỂN
% ==========================================
\chapter{KẾT LUẬN VÀ HƯỚNG PHÁT TRIỂN}
\section{Kết luận những đóng góp đạt được}
\section{Hạn chế của hệ thống}
\section{Định hướng phát triển tương lai}

\end{document}